\begin{abstract}
Long public radial-velocity archives can reveal low-mass planets, but archival signals
have a mixed record: many have been withdrawn as stellar activity, aliases or noise.
We measure how well common vetting tests separate planets from impostors on the real
sampling and noise of the HARPS-RVBank archive. We implement a ladder of interpretable
tests (false-alarm probability, alias comparison, phase and amplitude stationarity,
activity-indicator power and rotation harmonics) as tested open-source software and run
it blind on \nmSimCases\ simulated series built from \nmSimStars\ HARPS stars, with
Keplerian planets and evolving quasi-periodic activity injected into permuted real
noise. The false-alarm probability depends on how the jitter is fitted: held at its
planet-free value it is conservative by a factor of \nmFapFixFactorLo--\nmFapFixFactorHi,
and rescaled with the signal it is liberal by \nmFapLiberalPerm\ at 1\%. Many activity
impostors fall inside the broad lines of evolving rotation rather than on its exact
harmonics, more so the shorter the spot lifetime. No single test is decisive (best area under the ROC curve \nmAucInd, detection
statistics \nmAucFap); all rungs together reach \nmAucTree, with indicator power the most
important input. Of \nmBenchR\ published signals later refuted, \nmRefNotDet\ are not
significant in a blind search of the full archive; a rule whose thresholds were fixed before this benchmark was run (its features were
revised afterwards) rejects \nmRuleRejRef\ of the other \nmBRuleRefN\ and \nmBRuleUpRej\ of
\nmBRuleUpN\ significant upheld signals, counts that shift with the window allowed around
the claimed period. We also report the \nmDefN\ defects found in the
AI-assisted analysis that motivated this work, coded by two independent instances of the model: overclaims were
the largest category, few were numbers miscopied from code, and an independent reviewer
found most of them.
\end{abstract}

\begin{keywords}
{techniques: radial velocities -- methods: data analysis -- methods: statistical --
planets and satellites: detection -- stars: activity}
\end{keywords}

\maketitle
